\subsection{Gauges as Projective Fiber Transmittance}
  \label{sec:gauges-transmittance}

  Gauge interactions are reinterpreted as degrees of freedom of the
  admissible projective flow within the projection fiber~$\Pi$, encoding
  how locally admissible projections are coherently related despite
  non-injectivity.

  \begin{itemize}
    \item \textbf{Electromagnetism ($U(1)$):}
    Corresponds to the phase degree of freedom of the admissible projective
    flow along the fibers of the Hopf fibration of
    $\Pi \cong S^3$.
    This emergent geometric parameter reflects the freedom to locally
    re-identify projected descriptions without altering global relational
    consistency.
    The fine-structure constant~$\alpha$ characterizes the
    \emph{transmittance} of the projective flow through the fiber:
    the stability of admissible projective re-identifications along the
    fiber.

    \item \textbf{Weak Interaction ($SU(2)$):}
    Emerges from the rotational degrees of freedom of the $S^3$
    projection fiber itself.
    These encode non-abelian modes of projective transmittance,
    corresponding to shear-like distortions of admissible projections.
    The massive character of the $W^\pm$ and $Z^0$ bosons follows from a
    non-zero spectral gap associated with these modes, interpretable as
    spectral drag or torsion intrinsic to the fiber geometry.

    \item \textbf{Strong Interaction ($SU(3)$):}
    Associated with topological constraints arising from non-trivial
    winding structures of projected configurations.
    The derivation of an effective $SU(3)$-like symmetry from these
    constraints is an open problem: no mechanism selecting a threefold
    structure within the projection fiber is currently derived.
  \end{itemize}
